Mídias volumétricas, como nuvens de pontos e malhas, viabilizam experiências como realidade virtual e aumentada, mas depende de compressão com perdas, o que pode degradar a qualidade do conteúdo. Nesse contexto, medir a qualidade desses conteúdos tal como percebido pelo sistema visual humano se torna um problema crucial.
Esta monografia explora esse problema e desenvolve métricas de qualidade para esses mídias.

Este trabalho investiga três paradigmas complementares: (1) abordagens baseadas em projeção, que renderizam o conteúdo tridimensional em diferentes pontos de vista bidimensionais e aplicam métricas estabelecidas de qualidade de imagens, (2) abordagens baseadas em descritores que analisam a geometria e a a cor através de Análise de Componentes Principais, e (3) abordagens híbridas, que combinam tanto abordagens baseadas em projeções quanto descritores para capturar simultaneamente aparência e estrutura.

Tal investigação levou à produção de quatro grandes artefatos principais. (1) Uma métrica para nuvens de pontos baseada em projeção. Proposta neste trabalho, tal métrica atinge alta correlação com escores subjetivos. (2) ``PointPCA+RS'', uma reimplementação da métrica para nuvens de pontos chamada PointPCA+. Essa reimplementação é focada em alto desempenho computacional. (3) PointPCA++, uma métrica híbrida para nuvens de pontos que une tanto técnicas baseadas em projeções quanto descritores, obtendo maior robustez a diferentes distorções. (4) {MAPPE}, uma métrica para malhas inspirada na abordagem desenvolvida na PointPCA++, mas estendendo a capacidade de avaliação de qualidade para dados de malhas.

Em bases públicas, as métricas propostas alcançam o Coeficiente de Correlação de Postos de Spearman acima de 0,90 e, no caso de malhas, superam métodos do estado da arte anteriores com ganhos de mais de duas ordens de magnitude em termos de economia de tempo de processamento. Esses resultados indicam que métricas propostas para mídia volumétrica podem ser ao mesmo tempo perceptualmente fidedignas e viáveis para fluxos de trabalho imersivos.
